\section{Tutorial 12: Attention Visualization}

\subsection{Objectives}

This tutorial develops intuition for how attention behaves inside transformer models:

\begin{itemize}
    \item Understand attention weights as interaction patterns
    \item Interpret self-attention matrices visually
    \item Relate attention patterns to linguistic or structural roles
    \item Connect local and global dependencies to model behavior
\end{itemize}

\medskip

\noindent
\textbf{Focus.}  
Attention is not just a formula; it is a structured map of interactions between elements.

\subsection{Exercise 1: Attention Matrix as a Map of Interactions}

Consider a sequence of length $T$ with self-attention matrix:
\[
A = \text{softmax}\left(\frac{QK^\top}{\sqrt{d}}\right).
\]

\medskip

\noindent
\textbf{Tasks.}
\begin{enumerate}
    \item What are the dimensions of $A$?
    \item What does the entry $A_{ij}$ represent?
    \item Why does each row sum to $1$?
\end{enumerate}

\medskip

\noindent
\textbf{Solution.}

\begin{enumerate}
    \item
    \[
    A \in \mathbb{R}^{T \times T}.
    \]

    \item $A_{ij}$ is the attention weight from position $i$ to position $j$; it measures how much token $i$ attends to token $j$.

    \item Each row is the output of a softmax, so:
    \[
    \sum_{j=1}^T A_{ij} = 1.
    \]
\end{enumerate}

\medskip

\noindent
\textbf{Key takeaway.}  
Attention matrices are row-stochastic interaction maps.

\subsection{Exercise 2: Interpreting a Simple Attention Pattern}

Suppose a sentence has 5 tokens, and one attention head produces:

\[
A =
\begin{pmatrix}
0.7 & 0.2 & 0.1 & 0.0 & 0.0 \\
0.1 & 0.6 & 0.2 & 0.1 & 0.0 \\
0.0 & 0.2 & 0.5 & 0.2 & 0.1 \\
0.0 & 0.1 & 0.2 & 0.6 & 0.1 \\
0.0 & 0.0 & 0.1 & 0.2 & 0.7
\end{pmatrix}.
\]

\medskip

\noindent
\textbf{Tasks.}
\begin{enumerate}
    \item Describe the qualitative pattern.
    \item What kind of dependency does this head emphasize?
\end{enumerate}

\medskip

\noindent
\textbf{Solution.}

\begin{enumerate}
    \item Most of the mass is near the diagonal.

    \item This head emphasizes local context: each token mostly attends to itself and nearby tokens.
\end{enumerate}

\medskip

\noindent
\textbf{Insight.}  
Some attention heads behave like local filters, similar to convolutional receptive fields.

\subsection{Exercise 3: Global Attention Pattern}

Now consider:

\[
A =
\begin{pmatrix}
0.1 & 0.1 & 0.1 & 0.6 & 0.1 \\
0.0 & 0.2 & 0.1 & 0.6 & 0.1 \\
0.1 & 0.1 & 0.1 & 0.6 & 0.1 \\
0.1 & 0.1 & 0.1 & 0.5 & 0.2 \\
0.1 & 0.1 & 0.1 & 0.6 & 0.1
\end{pmatrix}.
\]

\medskip

\noindent
\textbf{Tasks.}
\begin{enumerate}
    \item What is the dominant attended position?
    \item What might this represent in language?
\end{enumerate}

\medskip

\noindent
\textbf{Solution.}

\begin{enumerate}
    \item Position 4 is dominant.

    \item This may correspond to a syntactically or semantically central token, such as a main verb, subject, or separator token.
\end{enumerate}

\medskip

\noindent
\textbf{Insight.}  
Attention can capture global structure, not only local neighborhoods.

\subsection{Exercise 4: Multi-Head Attention Interpretation}

Transformers use multiple attention heads:
\[
\text{MultiHead}(Q,K,V) = \text{Concat}(\text{head}_1, \dots, \text{head}_h)W_O.
\]

\medskip

\noindent
\textbf{Question.}  
Why is multi-head attention useful?

\medskip

\noindent
\textbf{Discussion.}

\begin{itemize}
    \item Different heads can specialize in different patterns
    \item One head may capture local context
    \item Another may capture long-range dependencies
    \item Another may track punctuation or structural boundaries
\end{itemize}

\medskip

\noindent
\textbf{Key takeaway.}  
Multi-head attention enables diverse relational representations.

\subsection{Exercise 5: Attention and Positional Structure}

Since self-attention alone is permutation-invariant, transformers add positional encodings.

\medskip

\noindent
\textbf{Tasks.}
\begin{enumerate}
    \item Why can attention alone not distinguish order?
    \item How does positional encoding fix this?
\end{enumerate}

\medskip

\noindent
\textbf{Solution.}

\begin{enumerate}
    \item Without positional information, the model only sees a set of vectors and their pairwise similarities.

    \item Positional encodings inject order information into token representations, so attention can depend on both content and position.
\end{enumerate}

\medskip

\noindent
\textbf{Insight.}  
Attention captures interactions; positional encoding tells the model where those interactions occur.

\subsection{Exercise 6: Comparison with RNNs and CNNs}

\textbf{Question.}  
How does attention differ from CNNs and RNNs in how it processes structure?

\medskip

\noindent
\textbf{Discussion.}

\begin{itemize}
    \item CNN:
    \begin{itemize}
        \item local interactions
        \item fixed receptive structure
    \end{itemize}

    \item RNN:
    \begin{itemize}
        \item sequential memory
        \item hidden-state bottleneck
    \end{itemize}

    \item Attention:
    \begin{itemize}
        \item global interactions
        \item data-dependent connectivity
    \end{itemize}
\end{itemize}

\medskip

\noindent
\textbf{Insight.}  
Attention replaces fixed or sequential structure with adaptive interaction patterns.

\subsection{Exercise 7: Attention as a Kernel Matrix}

Recall the matrix form:
\[
A = \text{softmax}\left(\frac{QK^\top}{\sqrt{d}}\right).
\]

\medskip

\noindent
\textbf{Task.}  
Explain why $QK^\top$ resembles a kernel or similarity matrix.

\medskip

\noindent
\textbf{Solution.}

\begin{itemize}
    \item Each entry is a dot-product similarity
    \item Similar to Gram matrices in kernel methods
    \item Difference: the representation itself is learned and input-dependent
\end{itemize}

\medskip

\noindent
\textbf{Insight.}  
Attention can be viewed as a learned, adaptive kernel operator.

\subsection{Exercise 8: Visualization Thought Experiment}

Suppose you plot attention matrices from different layers.

\medskip

\noindent
\textbf{Question.}  
What patterns might you expect from shallow vs deep layers?

\medskip

\noindent
\textbf{Discussion.}

\begin{itemize}
    \item Early layers:
    \begin{itemize}
        \item more local or lexical patterns
    \end{itemize}

    \item Later layers:
    \begin{itemize}
        \item more global, semantic, or task-specific interactions
    \end{itemize}
\end{itemize}

\medskip

\noindent
\textbf{Insight.}  
Attention patterns evolve hierarchically across depth.

\subsection{Exercise 9 (Discussion): Are Attention Weights Explanations?}

\textbf{Question.}  
Can attention maps be interpreted as explanations of model decisions?

\medskip

\noindent
\textbf{Discussion.}

\begin{itemize}
    \item They show interaction patterns
    \item But they do not fully explain the model’s computation
    \item Feedforward layers and residual paths also matter
\end{itemize}

\medskip

\noindent
\textbf{Insight.}  
Attention visualization is informative, but not a complete explanation.

\subsection{Summary}

\begin{itemize}
    \item Attention matrices visualize interactions between tokens
    \item Different heads can capture different structural roles
    \item Attention can express both local and global dependencies
    \item Transformers combine interaction modeling with positional information
\end{itemize}

\medskip

\noindent
\textbf{Preparation for the next part of the course.}  
We now move from architectures to generative modeling and large-scale AI systems, where attention and transformers play a central role.